\documentclass[12pt,letterpaper]{amsart}
\usepackage[T1]{fontenc}
\usepackage{mathptmx}
\usepackage{courier}
\usepackage[letterpaper,margin=1in]{geometry}
\usepackage{amsmath,amssymb,mathrsfs,enumitem}
\usepackage[hidelinks]{hyperref}
\setlength{\emergencystretch}{3em}
\DeclareMathOperator{\inv}{inv}

\title[Referee report: Shalika germs for $G_2$, Draft 6]{Referee report on ``On the transfer of Shalika germs for $G_2$'', Draft 6 (file \texttt{g2-shalika-germs-2026.tex})}
\author{Claude (Anthropic)}
\date{November 16, 2026}

\begin{document}
\maketitle

\section{Summary and overall assessment}

The manuscript is the November 2026 corrected and expanded edition of the 1994
manuscript. It proves the toric and patch results of \S\S3--4. It gives a
direct proof of Hypothesis $(\mathrm X_2)$ for the biquadratic tori
(Theorem~7.6H), conditional on Lemma~5.14 and the divisor proviso of
Lemma~0.8. For the cyclic cubic tori it reduces $(\mathrm X_3)$ to a sketched
principal-value computation (Claim~8.3H). On the whole the paper is honest
about its status: the conditional results are labeled, and the gaps are
recorded in Remarks~2026.$n$ and in the Open Problems.

I found no error that changes a stated result. The one compile error is fixed.
The main points are these:
\begin{enumerate}[label=(\roman*)]
\item The new remark on the repeated polar divisor
(\S\ref{sec:polar} below) has a consequence that the text does not draw out.
The tamely ramified line of (8.8) lies \emph{entirely} in the case
$q\equiv1\pmod 3$, so Claim~8.3H as stated covers directions $X$ for which the
sketch is not justified. I suggest a way around this that uses local constancy
of germs.
\item The Jacobian identity displayed in the proof of Theorem~7.6H(a)--(b) is
wrong. The conclusion (7.12) is right (I checked it independently).
\item A few local slips: an off-by-one count in Lemma~3.6, a wrong
cross-reference in \S4, an over-attribution to the divisor proviso in
Theorem~7.6H(d), and stale status statements together with an inconsistent
draft number.
\item On readability, the main problems are clashing symbols in
\S\S7--8 (some within a single proof), proofs in \S\S7--8 that consist mostly
of ``verified by script'', and a status summary repeated in seven places that
have drifted apart.
\end{enumerate}

\section{The compile error}

\texttt{\textbackslash mathscr} (used in Remark~2026, ``full regularity does not
exclude a repeated cubic polar divisor'', line~4888) is not defined by
\texttt{mathptmx} or \texttt{amssymb}. In the corrected file I added
\begin{verbatim}
\usepackage{mathrsfs}
\end{verbatim}
after \texttt{amsmath,amssymb,amsthm,amscd}. The displayed equation in that
remark was the only untagged numbered equation in the paper, so it printed as
``(1)''. I gave it \texttt{\textbackslash tag\{8.7i\}}, in keeping with the
paper's numbering. With these two changes the file compiles without errors,
undefined references or overfull boxes (73 pages). I made no other change to
the manuscript.

\section{What was checked independently}

These checks were made in exact rational arithmetic at random rational
points (scripts \texttt{s7num.py} and \texttt{s8num.py}, sent with this
report) or by hand. They are independent of the project's own scripts.

\begin{enumerate}[label=(C\arabic*)]
\item Table~4.3A. The values of $a$, $b$ and $\beta$ for all seven branches
follow from $u_E$ and $m_\omega=(7,-10,-6,1,1,1,1)$. For
$E_{\emptyset,\alpha,\delta}$ the last column is consistent with the exponent
vectors of the $x(\rho)$.
\item Lemma~3.5: the generator count is right. Lemma~3.6: the comparability
hypothesis holds for $G_2$ (but see \S\ref{sec:minor-correct}, item~1).
\item Proposition~7.6G(a), (b). The maps $s_\alpha^*$ and $s_\beta^*$ are
involutions and $(s_\beta^*s_\alpha^*)^6=1$. $(-1)^*$ acts on root values by
$-1$, and $(-1)^*\Phi_E|_{u_0}=\Phi_E|_{u_0}$.
\item Theorem~7.6H(b). The actions $s_\alpha^*:(U,V,\xi_1,\xi_2)\mapsto
(U^{-1},V,\xi_1,\xi_2)$, $(-1)^*\mapsto(U,V^{-1},\xi_2,\xi_1)$ and
$s_\zeta^*\mapsto(U^{-1},V^{-1},\xi_2,\xi_1)$ hold exactly. The pullback
(7.12) holds with sign $-1$.
\item Theorem~7.6H(c). The factorization of $\Phi_E|_{u_0}$ and the identity
$\omega_{K/L}(\Phi_E|_{u_0})=(t^2-d_\alpha d_\zeta,d_\alpha)_F$ check out by hand.
\item Theorem~7.6H(d). I recomputed $I_t$, $I_s$ and $I_\xi$ by hand in the
three valuation cases, using the analytic-continuation regularization
($\mathrm{PV}\!\int_{|t|>1}dt=-1$, $\mathrm{PV}\!\int_{\mathcal O}|y|^{-2}dy=-q^{-1}$),
and obtained the values stated. The products give (7.13) in all three
cases.
\item Example~7.6B(b). For $p=5$ the Legendre symbol is $-1$ for all unit
residues. For $p=7,11,\dots,29$ it takes both values.
\item Proposition~8.3F(a). In the intrinsic coordinates (8.7d), the
stated formulas for $s_\alpha^*$ and $s_\beta^*$ hold exactly, and so do
$s_\alpha^*(L_X)=-L_X/D_{s_\alpha}^2$ and $s_\beta^*(L_X)=L_X/r_\gamma^3$. For
this I took the geometric actions to be those of Proposition~2.14A, with the
root values transformed as in Proposition~7.6G, extended by
$s_\alpha^*(\eta)=\eta+\alpha z(\alpha)$.
\item Claim~8.3H. I recomputed $I_{\mathbb P^1(K)}=-2(q+1)q^{-2}$ by hand. The
measure reduces to $dc_0\,dc_1\,dc_2/|N_{K/F}(c_1+c_2\varphi)|$, and I
regularized with $|N|^{s}$. The value is the same in the unramified case
($-(q+1)q^{-2}$ from $\mathcal O_K$ and the same again from the complement) and
in the tame case ($-2q^{-1}$ and $-2q^{-2}$). The passage from $B_1'(1)B_2(1)$
to (8.8) is arithmetically consistent.
\item The algebra of the repeated-polar remark is correct:
$M^2+3\Delta^2=12(3a^2+3ab+b^2)=6\operatorname{Tr}_{K/F}(\alpha(X)^2)$. Every
nonzero $X\in\mathfrak t(F)$ is regular. The zero set is a pair of eigenlines
exactly when $\mu_3\subset F$. The $\mathbb Q_7$ example is valid: $2$ is not
a cube mod~$7$, so $x^3-2$ gives the unramified cubic extension.
\end{enumerate}

\section{Correctness}

\subsection{The repeated-polar remark and Claim~8.3H}\label{sec:polar}

\begin{enumerate}
\item \emph{Scope.} A tamely ramified cyclic cubic extension exists only when
$\mu_3$ lies in the residue field, that is, $q\equiv1\pmod3$. So the second
line of (8.8) is \emph{always} in the repeated-polar case. The first
line is in it whenever $q\equiv1\pmod 3$. As stated, Claim~8.3H covers these
$X$, but by the remark's own account the one-parameter principal-value
argument is not justified there. The Claim (or Remark~8.status(b)) should say
so plainly.
\item \emph{A suggested repair.} Under the Claim's own hypotheses (Lemma~5.14
and the proviso), Corollary~8.3B and (8.6) identify
$J_3(X)/|\alpha\gamma\delta(X)|$, up to $X$-independent factors, with the
normalized $3$-regular $\kappa$-germ. Shalika germs are locally constant on
the regular set near $1$. The set $\mathscr D(X)=0$ is a union of two lines,
and every point of it is a limit of regular points off it. So it would be
enough to prove the identity for $\mathscr D(X)\neq0$ and extend by local
constancy. This avoids any regularization at the exceptional pole. Please
check that the identification really involves only factors that are
continuous in $X$ (the $\kappa(\inv)$ factor and $\Delta_0$ are constant here
by Proposition~8.3A).
\item \emph{Undefined terms.} The remark uses ``the quadratic polar
polynomial'', ``the exceptional quadric cone'', ``the associated
exceptional-divisor resonance'' and ``the auxiliary analysis''. None of these
is defined in the paper, and none is tied to the objects of Claim~8.3H
($\ell_1$, $t_*(X)$, the coalescence at $U=-3$). Also, the cone
$\mathcal C_K$ of Proposition~8.3G is called smooth and split, not
exceptional. Give the polynomial explicitly and say where its discriminant
enters the sketch, or say only what the external note
\texttt{notes/cubic-regularity-discriminant-2026-11-15.md} establishes.
\item \emph{Placement.} The remark sits between Proposition~8.3G and
Claim~8.3H but is about the proof of the Claim. It would read better directly
after the Claim's sketch, or folded into Remark~8.status(b).
\item \emph{A detail.} The argument that coinciding conjugates force $X=0$
uses only trace zero and characteristic $0$; ``$p\geq5$'' is not needed there.
The phrase ``This condition is \emph{not} implied to be nonzero by regularity''
should be reworded, e.g.\ ``Regularity of $X$ does not imply
$\mathscr D(X)\neq0$.''
\end{enumerate}

\subsection{Theorem~7.6H, proof of (a)--(b)}

The displayed identity
\[
\det\frac{\partial(U,\log V)}{\partial(w(\gamma),w(\delta))}
=\frac{\alpha\zeta H}{F_0^3F_1F_2}
\]
is false. The correct value is
\[
\det\frac{\partial(U,\log V)}{\partial(w(\gamma),w(\delta))}
=\frac{\alpha\zeta H}{F_0^2\,G},
\qquad\text{equivalently}\qquad
\det\frac{\partial(U,V)}{\partial(w(\gamma),w(\delta))}
=\frac{\alpha\zeta H}{F_0^3F_1F_2}.
\]
So the printed right side is the Jacobian of $(U,V)$, not of $(U,\log V)$. The
middle expression in the display, with the factor
$(U-1)^2/(F_0^2/(\alpha^2\zeta^2w(\delta)^2))$, does not simplify to either
value: by $U-1=\alpha\zeta w(\delta)/F_0$ it equals
$\alpha^4\zeta^4w(\delta)^4/F_0^4$. With the corrected determinant, and
$d\xi_1\wedge d\xi_2/(\xi_1-\xi_2)^2\equiv (G/(Hz(\beta)^2))\,d\xi\wedge dz(\beta)$
modulo $dw(\gamma),dw(\delta)$, one recovers (7.12) exactly. I
confirmed this (C4). The symbols $L_1$, $L_2$ and $Q_0$ in ``$V=L_2/(L_1Q_0)$''
are not defined. $Q_0$ is later defined differently in Proposition~8.3G.

\subsection{Theorem~7.6H, proof of (d)}

The proof says: ``Under the divisor proviso of Lemma~0.8, no other branch $E'$
of $\phi^{-1}(1)$ on a smooth resolution has $\beta(E')=3$.'' The proviso
excludes only branches with the same $(\theta,\beta)$ \emph{whose $F$-points
map onto $\mathcal O^{st}(F)$}. What the argument needs is that no branch with
$(\theta,\beta)=(1,3)$ \emph{meets} the strict transform of
$E_{\emptyset,\alpha,\delta}$, and that is $e(1,3)=1$, i.e.\ Lemma~5.14. (A
branch with $\theta\neq1$ gives a factor regular at $s=0$ anyway.) Please
reword so that each input is cited for what it supplies.

\subsection{Smaller correctness points}\label{sec:minor-correct}

\begin{enumerate}
\item \emph{Lemma~3.6, proof.} The count
``$\ell+|\{\rho\ge\sigma\}|+|\{\rho<\sigma\}|-\ell=p+1$'' evaluates to $p$. The
$w$-generators are the non-simple $\rho\le\sigma$, which number
$|\{\rho<\sigma\}|+1-\ell$, because $\sigma$ itself counts and all simple roots
lie below a non-simple $\sigma$. The count should read
$\ell+|\{\rho\ge\sigma\}|+(|\{\rho<\sigma\}|+1-\ell)=p+1$. The hypothesis that
all simple roots are $<\sigma$ should be stated. It holds for $G_2$.
\item \emph{\S4, after Table~4.3A.} ``By Igusa's theorem (property (e) in
\S0)'': property~(e) of \S0 is faithful flatness of $\pi|_E$. The intended
reference is the expansion (0.3)/(5.5), or
Lemma~5.7.
\item \emph{Remark~5.14} says the argument ``remains incomplete in three ways''
and then lists four items. Item~(iv) is a removed sentence, not an
incompleteness. Say ``three ways (i)--(iii)'' and make (iv) a separate
sentence.
\item \emph{Remark~2026 on the main gap, last paragraph.} It says that for
$r=3$ only Propositions~8.3A--8.3E are available and that the evaluation of
the $6$-dimensional integral ``remain[s] open''. That predates
Propositions~8.3F--8.3G and Claim~8.3H. Update it to match Remark~8.status.
\item \emph{Draft number.} The title, date line and header comment say
Draft~6. But the text records changes made \emph{after} Draft~6: Claim~8.3H
says ``Draft~6 stated $\pm\frac{2}{27}$''; the Acknowledgments say
``Claude refereed Draft~6, corrected the residue''; and the repeated-polar
remark cites a note dated 2026-11-15. By the style guide this should be
Draft~7. The Acknowledgments should then say who wrote the repeated-polar
remark and its verification script, and Open Problem~(1) should say
``as of Draft~7''.
\item \emph{Proposition~8.3F(b):} ``commutes elementwise with $\Omega(G_2)$''
is unclear. $SL_2(\alpha)$ is not normalized pointwise by the Weyl group in
any obvious sense. Presumably what is meant is that the $SL_2(\alpha)$-action
commutes with the geometric $\Omega(G_2)$-action on the fiber. Please say so,
or give the reason.
\item \emph{Claim~8.3H} factors $J_3(X)$ as a product of separately
regularized integrals ($I_{\mathbb P^1(K)}$ is the regularized volume of a
$PGL_2(F)$-torsor, which has infinite measure). That a regularized integral
over a product equals the product of the regularized factors is part of what
Remark~8.status(b) must justify, and could be named there explicitly. The
value of $I_{\mathbb P^1(K)}$ itself is correct (C9).
\end{enumerate}

\section{Readability}

\subsection{Clashing symbols that matter}

The paper already acknowledges the root/exponent clashes inherited from 1994
(``$\beta$'', ``$\gamma,\delta,\epsilon,\zeta$''). The following are new, and
several occur inside a single proof or proposition:
\begin{itemize}
\item $L$: the field $F(\sqrt{d_\alpha d_\zeta})$, and also
$L:=x(\delta)z(\beta)/w(\delta)$ in the proof of Proposition~7.6E. That
proof uses $\omega_{K/L}$ a few lines later. This is the worst case; rename
the function (e.g.\ $\ell_E$).
\item $u_0$: the basepoint $\theta_\delta(1)$ in \S7, the basepoint
$\theta_\zeta(1)$ in \S8, and the cone coordinate
$u_0:=1+2(j_3'+j_3'')/j_2^2$ in Proposition~8.3G, whose part~(b) uses both
meanings (``on $E_{\emptyset,\beta,\zeta}(u_0)(F)$''). Rename the cone
coordinates (e.g.\ $c_0,c_1,c_2$).
\item $t_*(X)$ in Claim~8.3H (the root of $\ell_1$) versus the cocycle $t_*$.
\item $G$, $H$ (polynomials in Propositions~7.6E--7.6G) versus the group and
the endoscopic group; $F_0,F_1,F_2$ versus the field $F$; $P(Y)$, $P_1$,
$P_2$ versus the parabolic $P$ and the weight lattice $P$; $\Delta$ and $M$
in the repeated-polar remark versus the simple roots, the transfer factor, and
$M$ in Lemma~5.7 and \S A.7; $\delta_L$, $\delta_K$ versus the root $\delta$;
$Q$, $Q_0$, $Q_1$, $Q_2$, $\mathcal Q_3$ and the root lattice $Q$.
\end{itemize}
A short ``Notation for \S\S7--8'' paragraph would help, or an extension of the
Notation table.

\subsection{Proofs that only cite a script}

The proofs of Propositions~7.6E, 7.6G, 8.3D, 8.3E, 8.3F and 8.3G consist mostly
of ``exact polynomial composition'' plus a script name. As written, a reader
cannot check them without running the code. For each, give the key
intermediate formulas: the leading Laurent terms of the $n_\rho$ in 7.6E;
$q_1$ and $Q_1$ explicitly in 8.3C--8.3D; the cone parametrization and the
Hilbert-90 element in 8.3G. If space is a concern, put them in an appendix.
The Claim~8.3H sketch already admits that $P_0$, $\ell_1$, $L_0$, $L_1$,
$N_K$, $s(U,t)$, $C(X)$ and $c_{\mathrm{alg}}$ are undefined; this belongs at
the top of Open Problem~(1)(b).

\subsection{Structure}

\begin{enumerate}
\item \emph{Forward-referencing proofs.} The ``Proof of Theorem~7.2'' ends with
$\square$ after saying the needed statements ``are proved \dots\ below''. The
same pattern occurs for Theorem~8.3. Either move the proof heading after
Theorem~7.6H (respectively Claim~8.3H), or make Theorem~7.6H the main theorem
of \S7 and derive Theorem~7.2 from it in two lines.
\item \emph{Status is stated seven times.} It appears in the abstract, the
``Status of the results'' list in \S0, the main-gap remark, Remark~7.status,
Remark~8.status, the Summary, and the Open Problems, and these have drifted
(\S\ref{sec:minor-correct}, item~4). Keep one authoritative list (the Summary)
and replace the others by a sentence and a pointer.
\item \emph{Density of Remarks~2026.n.} There are about sixty. Many only
record a 1994 misprint (e.g.\ the remarks after (1.18), the
\S3 patch formulas, the \S4 forms, (2.6), the regular branch).
The header states that clear-cut changes go to \texttt{CHANGES.md}. Moving the
pure-misprint remarks there, or into an appendix ``Corrections to the 1994
text'', would make the mathematics much easier to follow. Keep in the text
the remarks that affect the logic (transfer known, 0.8, 5.14, exceptional
divisors, main gap, 6.1, 8.roots, 7.status, 8.status, repeated polar).
\item \emph{Revision history inside the mathematics} (``Draft~4 consulted
\dots'', ``Draft~5 also supplied \dots'', ``Draft~6 stated $\pm\frac{2}{27}$'')
belongs in \texttt{revision\_log.md}, per the style guide.
\item \emph{Abstract.} At about 30 lines it is long and cites five
proposition numbers. Shorten it to the AI-disclosure sentence, what is proved,
what is conditional, and the fact that the matching is now known by
\cite{GH16}; the details belong in \S0.
\item \emph{Labels.} Several labels do not match the printed numbers
(\texttt{res:7.6} is Proposition~7.6A, \texttt{res:7.7} is Example~7.6B,
\texttt{res:7.8} is Corollary~7.6C). This is invisible in the PDF but a trap
for whoever edits next.
\end{enumerate}

\section{Recommendation}

The corrections in \S4 are all local. The repeated-polar issue should be
handled by narrowing Claim~8.3H or by the local-constancy argument suggested
in \S\ref{sec:polar}, item~2. On readability, the most valuable changes are
fixing the $L$ and $u_0$ clashes, consolidating the status statements, and
moving the misprint remarks to \texttt{CHANGES.md}. After these changes the
paper would be a clear record of what Igusa theory does and does not yet
prove for $G_2$.

\begin{thebibliography}{GH16}
\bibitem[GH16]{GH16} J.~Gordon and T.~Hales, Endoscopic transfer of orbital
integrals in large residual characteristic, Amer. J. Math. \textbf{138}
(2016), 109--148.
\end{thebibliography}

\end{document}
